\documentclass[10pt,a4paper]{amsart}
\usepackage[margin=19mm]{geometry}
\usepackage{amsmath,amssymb,mathtools}
\usepackage[scr=rsfs,cal=euler]{mathalfa}
\usepackage{xfrac}
\usepackage[unicode,hidelinks,bookmarks=false]{hyperref}
\newcommand{\ie}{i.e.\ }
\newcommand{\cf}{cf.\ }
\newcommand{\resp}{respectively\ }
\newcommand{\Subsection}{\subsection}
\providecommand{\nobreakdash}{\nobreak\hbox{-}}
\setlength{\emergencystretch}{2em}
\multlinegap0pt
\newcommand{\II}{{\rm II}}
\usepackage{amssymb}
\usepackage{enumerate}
\newtheorem{lemma}{Lemma}
\newtheorem{theorem}{Theorem}
\newcommand{\III}{{\rm III}}
\newcommand{\R}{\mathbb{R}}
\newcommand{\cR}{\mathcal{R}}
\newcommand{\cT}{\mathcal{T}}
\newcommand{\rL}{\mathord{\text{\rm L}}}
\title{Uniqueness of almost periodic outer flows on the hyperfinite type $\II_1$ factor}
\author{Cyril Houdayer, Amine Marrakchi}
\date{Source: arXiv:2605.02781v2 (CC BY 4.0)}
\begin{document}
\maketitle

\noindent\emph{Source and licence.} Uniqueness of almost periodic outer flows on the hyperfinite type $\II_1$ factor. Version arXiv:2605.02781v2 (CC BY 4.0), distributed by arXiv under the Creative Commons Attribution 4.0 International licence, http://creativecommons.org/licenses/by/4.0/, as recorded in the arXiv licence metadata for this exact version and shown on the abstract page https://arxiv.org/abs/2605.02781v2. Full licence notice: \texttt{licenses/arxiv-2605.02781-CC-BY-4.0.txt}.\par\smallskip

\noindent\textit{Editorial scope.} Counted unit: the article's theorem subgroup (source0.tex lines 551--561). The article's complete original proof of it is delivered with it (source0.tex lines 585--613, one \texttt{proof} environment of 29 lines, which the article places away from the statement). No mathematics is written, repaired or re-derived: the statement and the proof are the article's own lines, in the article's own notation, and no retained line is substituted or re-flowed.  Retained uncounted as the source-local context the proof needs: lines 502--504; lines 566--576.  Nothing was omitted from any retained span, so the package carries no omission record.  No retained line contains a citation, so the wrapper carries no bibliography and no bibliography entry.  No retained line contains a figure environment, an image-inclusion command or drawing code, so no image is delivered and the package declares no asset dependency.  Editorial formatting only, in addition: an AMS \texttt{amsart} preamble that restates the article's own declarations and macros used by the retained lines, namely the environments \texttt{lemma}, \texttt{theorem} on separate counters, together with the standard packages those lines need.  The retained environments print in this document as Lemma 1, Theorem 1 in order of appearance, whereas the article prints the same items with the numbering of its own full document.  The complete native source of the article as distributed in the arXiv e-print for this exact version is bundled unchanged as \texttt{sources/199746/source0.tex}.
\begin{lemma} \label{dense range induced equivalence relation}
    Let $\cR$ be an ergodic nonsingular equivalence relation on $(X,\nu)$. Let $c : \cR \rightarrow G$ be a cocycle with values in a second countable locally compact group $G$. Let $\rho : G \curvearrowright (Y,\eta)$ be a nonsingular action. If $c$ is ergodic, then $\rL^\infty(X \times Y, \nu \otimes \eta)^{\cR_{\rho,c}} = \rL^\infty(Y)^{\rho}$.
\end{lemma}

\begin{theorem} \label{model ergodic kernel dense subgroup measure preserving case}
    Let $G$ be a second countable locally compact group. Then there exists an ergodic probability measure preserving equivalence relation $\cR$ with an ergodic cocycle $c : \cR \rightarrow G$.

    Moreover, the following holds :
    \begin{enumerate}
        \item If $G$ is amenable, we can take $\cR$ to be amenable.
        \item We can take $c$ such that $\ker(c)$ is ergodic.
        \item If $\Lambda < G$ is a dense countable subgroup, we can take $c$ with values in $\Lambda$.
    \end{enumerate}
    Any two of these conditions can be satisfied simultaneously. If $\Lambda$ is amenable, the three of them can be satisfied simultaneously.
\end{theorem} 

\begin{theorem} \label{model ergodic kernel dense subgroup}
    Let $G$ be a second countable locally compact group. Let $\rho : G \times \R^*_+ \curvearrowright (Y,\eta)$ be any ergodic nonsingular action. Then there exists an ergodic nonsingular equivalence relation $\cR$ on a probability space $(X,\nu)$ and a cocycle $c : \cR \rightarrow G$ such that the Mackey range of $c \times \omega : \cR \rightarrow G \times \R^*_+$ is isomorphic to $\rho$, where $\omega : \cR \rightarrow \R^*_+$ is the Radon--Nikodym cocycle.

    Moreover, the following holds :
    \begin{enumerate}
        \item If $G$ is amenable, we can take $\cR$ to be amenable.
        \item If $\rho|_{\R^*_+}$ is ergodic, we can take $c$ such that $\ker(c)$ is ergodic.
        \item If $\Lambda < G$ is a dense countable subgroup, we can take $c$ with values in $\Lambda$.
    \end{enumerate}
    Any two of these conditions can be satisfied simultaneously. If $\Lambda$ is amenable, the three of them can be satisfied simultaneously.
\end{theorem}

\begin{proof}[Proof of Theorem \ref{model ergodic kernel dense subgroup}]
    Let $\rho : G \times \R^*_+ \curvearrowright (Y,\eta)$ be an ergodic nonsingular action. Thanks to Theorem \ref{model ergodic kernel dense subgroup measure preserving case}, we can take a measure preserving equivalence relation $\cR_0$ on a probability space $(X_0,\nu_0)$ and an ergodic cocycle $c_0 : \cR_0 \rightarrow G$. Let $\cT$ be an ergodic amenable nonsingular equivalence relation on a probability space $(T,\zeta)$ such that the Radon--Nikodym cocycle $\delta : \cT \rightarrow \R^*_+$ is ergodic, i.e.\ $\cT$ is of type $\III_1$. 
    
    On $(X,\nu)=(X_0 \times T \times Y, \nu_0 \otimes \zeta \otimes \eta)$ consider the induced nonsingular equivalence relation $\cR=(\cR_0 \times \cT)_{\rho,c_0 \times \delta}$. Define the cocycle 
    $$c : \cR \ni (x,t,y,x',t',y') \mapsto c_0(x,x') \in G.$$

    Let us compute the Mackey range of $c \times \omega$ where $\omega$ is the Radon--Nikodym cocycle of $\cR$ with respect to $\nu$.
    
     Let $\sigma : G \times \R^*_+ \times Y \rightarrow \R^*_+$ be the Radon--Nikodym cocycle of $\rho$. Then the cocycle $\omega : \cR \rightarrow \R^*_+$ is given by  $$ \omega(x,t,y,x',t',y') = \delta(t,t') \sigma(c_0(x,x'),\delta(t,t'),y').$$
     

    The product cocycle $c \times \omega : \cR \rightarrow G \times \R^*_+$ is given by
    $$ (c \times \omega)(x,t,y,x',t',y') =\left( c_0(x,x'), \delta(t,t') \sigma(c_0(x,x'),\delta(t,t'),y') \right).$$

     
     
    We want to compute the algebra $$\rL^\infty(X_0 \times T \times Y \times G \times \R^*_+)^{\cR_{c \times \omega}}$$ of all invariant functions for the skew product equivalence relation $\cR_{c \times \omega}$. 
    
    Consider the action $\kappa :  G \times \R^*_+ \curvearrowright Y \times G \times \R^*_+$ given by $$(g,\lambda) \cdot (y,h,s) = (\rho(g,\lambda) y, g h, \lambda \sigma(g,\lambda,y) s).$$
    Then we have the equality $$\cR_{c \times \omega} = (\cR_0 \times \cT)_{\kappa,c_0 \times \delta}.$$ 
    Since $c_0 \times \delta : \cR_0 \times \cT \rightarrow G \times \R^*_+$ is ergodic, we conclude by Lemma \ref{dense range induced equivalence relation} that $$\rL^\infty(X_0 \times T \times Y \times G \times \R^*_+)^{\cR_{c \times \omega}}=\rL^\infty(Y \times G \times \R^*_+)^{\kappa}.$$

     We can regard these $\kappa$-invariant functions in $\rL^\infty(Y \times G \times \R^*_+)^{\kappa}$ as $G$-equivariant functions in $\rL^\infty(G, \rL^\infty(Y)^{(1,\alpha)})$ where $\alpha=\rho|_{\R^*_+}$ and $G$ acts on $\rL^\infty(Y)^{(1,\alpha)}$ by the restriction of the Maharam extension of $\rho|_G$. The Mackey range of $c \times \omega$, which is the restriction of the right translation action of $G \times \R^*_+$ on $\rL^\infty(Y \times G \times \R^*_+)^{\kappa}$, is then identified with $\rho^{(1)}$. Therefore, if we replace $\rho$ by $\rho^{(-1)}$ in our construction, we obtain that the Mackey range of $c \times \omega$ is precisely $(\rho^{(-1)})^{(1)} \cong \rho$.

    Observe that if  $\cR_0$ is amenable, then $\cR$ is also amenable. If $c_0$ takes its values in a dense countable subgroup $\Lambda < G$, then so does $c$. 

    Suppose that $\alpha=\rho|_{\R^*_+}$ is ergodic and choose $c_0$ such that $\ker(c_0)$ is ergodic. Observe that $\ker(c)$ is the induced equivalence relation $(\ker(c_0) \times \cT)_{\alpha,1 \times \delta}$.
     Since the cocycle $1 \times \delta  : \ker(c_0) \times \cT \rightarrow \R^*_+$ is ergodic then, thanks to Lemma \ref{dense range induced equivalence relation}, we conclude that $\ker(c)$ is ergodic.
\end{proof}

\end{document}
